\chapter{Q: Die Quelle der Selbstorganisation}

\section{Einleitung}

Die Axiome der IWT definieren die Struktur der Information – aber sie erklären nicht, warum aus reiner Information stabile Strukturen entstehen.

Die Antwort ist \( Q \).

\( Q \) ist das Prinzip der Selbstorganisation. Es ist die Kraft, die die Information so ordnet, dass sie stabil bleibt. \( Q \) ist kein externes Postulat – es folgt zwingend aus den Axiomen der IWT.

Dieses Kapitel zeigt, wie \( Q \) aus den Axiomen hergeleitet wird, und welche Rolle es in der Theorie spielt. Darüber hinaus wird die Interpretation von \( I \) als Energie, Masse und Ladung dargestellt – und die Definition der Ladung als diskrete Divergenz der Phase aus physikalischer Plausibilität und Konsistenz mit der emergierten Physik (insbesondere der Weber-Elektrodynamik) begründet.

\section{Q ist kein Axiom}

\( Q \) ist nicht in den Axiomen der IWT enthalten.

Die Axiome definieren:
\begin{itemize}
\item die Existenz von Information \( I_k^{(n)} \),
\item ihre Erhaltung \( \sum |I|^2 = \text{const} \),
\item ihre Dynamik \( \delta \mathcal{L}_d = 0 \).
\end{itemize}

\( Q \) ist keine zusätzliche Annahme. \( Q \) ist eine Konsequenz.

Es ist das, was passiert, wenn Information unter der Bedingung der Erhaltung und des Variationsprinzips sich selbst organisiert.

\section{Die Interpretation von I als Energie, Masse und Ladung}

Die fundamentale Größe der IWT ist die Information \( I_k^{(n)} \).

Sie ist die einzige ontologische Größe der Theorie.

Alle physikalischen Größen – Energie, Masse, Ladung – sind Interpretationen von \( I \). Sie sind nicht fundamental, sondern emergente Aspekte derselben Informationsdynamik.

\subsection{Energie}

Energie ist der Informationsfluss:

\[
E^{(n)} = \sum_k \left[ \frac{\partial \mathcal{L}_d}{\partial (\Delta_t I_k^{(n)})} \Delta_t I_k^{(n)} - \mathcal{L}_d \right]
\]

Sie ist die zum Zeitschritt \( T \) konjugierte Größe. Sie ist keine fundamentale Größe – sie ist eine Eigenschaft der Informationsdynamik.

\subsection{Masse}

Masse ist der Widerstand gegen Informationsänderung:

\[
m_k^{(n)} = \delta \sum_{l \in \mathcal{N}(k)} |I_k^{(n)} - I_l^{(n)}|^2
\]

Sie ist ein Maß für die lokale Struktur der Information. Je stärker die Information an einem Knoten von seiner Umgebung abweicht, desto größer ist die effektive Masse.

\subsection{Ladung}

Die Definition der Ladung in der IWT folgt aus physikalischer Plausibilität und der Konsistenz mit der emergierten Physik – insbesondere mit der Weber-Elektrodynamik (WED) und der Rotverschiebung.

Die Information \( I \) ist komplex und besitzt zwei Aspekte:
\begin{itemize}
\item die \textbf{Amplitude} \( |I| \), welche die Intensität, Dichte und Masse beschreibt,
\item die \textbf{Phase} \( \phi = \arg(I) \), welche die Kohärenz, Ausbreitung und das Feld beschreibt.
\end{itemize}

Die Masse ist die Änderung der Amplitude:
\[
m = \nabla |I|
\]

Die Ladung ist die Änderung der Phase – aber nicht der einfache Gradient, sondern die Divergenz des Phasengradienten:

\[
q = \nabla^2 \phi
\]

Im diskreten Fall der IWT bedeutet dies:

\[
\boxed{q_k^{(n)} = \sum_{l \in \mathcal{N}(k)} \left( \phi_k^{(n)} - \phi_l^{(n)} \right)}
\]

Begründung:

\begin{enumerate}
\item \textbf{Masse und Ladung sind komplementär:} Die Masse ist der Widerstand gegen Änderung der Amplitude (Materie). Die Ladung ist die Divergenz der Phasenänderung (Feld). Sie sind zwei Aspekte derselben Information – aber sie sind nicht identisch.

\item \textbf{Ladung ist die Quelle des Feldes:} In der WED ist die Ladung die Quelle des elektrischen Potentials. Die WED ist eichinvariant – sie ist invariant unter der Transformation \( \phi \to \phi + \chi \), wobei \( \chi \) eine beliebige Funktion ist. Die Ladung muss daher eichinvariant sein. Der einfache Phasengradient \( \nabla \phi \) ist nicht eichinvariant – er ändert sich unter \( \phi \to \phi + \chi \). Seine Divergenz \( \nabla^2 \phi \) hingegen ist eichinvariant, weil sie unter der Transformation unverändert bleibt. Im diskreten Fall ist die Summe der Phasendifferenzen die natürliche diskrete Divergenz.

\item \textbf{Die Rotverschiebung ist eine Phasenänderung:} Photonen verlieren Energie durch die Änderung ihrer Phase. Die Rotverschiebung ist die Energiesenke, die das Universum stabil hält. Sie ist mit der Phase und damit mit der Ladung gekoppelt.

\item \textbf{Die Fluktuation aus der diskreten Zeit verändert die Ladung nicht:} Die Vakuumfluktuation (aus Anhang O) erzeugt Struktur, aber sie erzeugt oder vernichtet keine Ladung. Die Ladung ist eine Erhaltungsgröße.
\end{enumerate}

\subsection{Die Einheit}

Energie, Masse und Ladung sind drei Interpretationen derselben Größe: der Information \( I \).

\[
\begin{aligned}
\text{Energie} &= f_{\text{Energie}}(I) \\
\text{Masse} &= f_{\text{Masse}}(I) \\
\text{Ladung} &= f_{\text{Ladung}}(I)
\end{aligned}
\]

Alle physikalischen Größen sind \( I \).

\( I \) ist die einzige fundamentale Größe.

\section{Die Herleitung von Q}

Die Herleitung von \( Q \) folgt aus zwei Prinzipien:

\begin{enumerate}
\item \textbf{Informationserhaltung:} \( \sum |I|^2 = \text{const} \)
\item \textbf{Minimale Spannung:} Die Information wird so organisiert, dass die Änderung zwischen benachbarten Knoten minimal ist.
\end{enumerate}

\subsection{Der Ansatz}

Wir betrachten ein Funktional, das die Spannung im Netzwerk misst:

\[
\mathcal{S} = \sum_k |\nabla I_k|^2
\]

Dabei ist \( \nabla I_k \) der diskrete Gradient der Information am Knoten \( k \).

Die Informationserhaltung wird als Nebenbedingung eingeführt:

\[
\sum_k |I_k|^2 = \text{const}
\]

\subsection{Die Lagrange-Multiplikator-Methode}

Wir minimieren die Spannung unter der Nebenbedingung der Informationserhaltung:

\[
\mathcal{L} = \sum_k |\nabla I_k|^2 + \lambda \left( \sum_k |I_k|^2 - \text{const} \right)
\]

Die Variation nach \( I_k \) ergibt:

\[
\delta \mathcal{L} = 0
\]

Daraus folgt die Euler-Lagrange-Gleichung:

\[
\Delta^2 I_k = \lambda \cdot I_k
\]

Dabei ist \( \Delta^2 \) der diskrete Laplace-Operator.

\subsection{Die Lösung}

Die Lösung dieser Gleichung hat die Form:

\[
I_k = \sqrt{\rho_k} \cdot e^{i \phi_k}
\]

wobei \( \rho_k = |I_k|^2 \) die Wahrscheinlichkeitsdichte und \( \phi_k \) die Phase der Information ist.

Einsetzen in die Euler-Lagrange-Gleichung und Trennung von Real- und Imaginärteil ergibt zwei Gleichungen.

Die Imaginärteil-Gleichung liefert die Kontinuitätsgleichung:

\[
\frac{\partial \rho}{\partial t} + \nabla \cdot (\rho \vec{v}) = 0
\]

Die Realteil-Gleichung liefert die Hamilton-Jacobi-Gleichung mit einem zusätzlichen Potential:

\[
\frac{\partial S}{\partial t} + \frac{|\nabla S|^2}{2m} + V + Q = 0
\]

wobei:

\[
Q = -\frac{\hbar^2}{2m} \frac{\nabla^2 \sqrt{\rho}}{\sqrt{\rho}}
\]

\subsection{Die Identifikation von Q}

Das Potential \( Q \) ist das Bohm-Potential der De-Broglie-Bohm-Theorie.

Es ist die Quelle der Selbstorganisation.

Es ist die Kraft, die die Information bindet und stabilisiert.

Es ist \( Q \).

\section{Q als Quelle}

\( Q \) ist die Quelle von allem:

\begin{enumerate}
\item \textbf{DBT:} \( Q \) ist das Bohm-Potential der De-Broglie-Bohm-Theorie.
\item \textbf{WG:} \( \vec{F}_{\text{WG}} = -\nabla Q \) – die Weber-Gravitation ist der Gradient von \( Q \).
\item \textbf{WED:} \( \vec{F}_{\text{WED}} = -\nabla Q \) – die Weber-Elektrodynamik ist der Gradient von \( Q \).
\end{enumerate}

Alles kommt aus \( Q \).

\section{Q als Selbstorganisation}

\( Q \) ist das Prinzip der Selbstorganisation.

Die Information organisiert sich selbst, weil sie unter der Bedingung der Erhaltung und der minimalen Spannung stabil sein muss.

\( Q \) ist der Name dieser Organisation.

\( Q \) ist die Quelle.

\section{Zusammenfassung}

\begin{itemize}
\item \( Q \) ist kein Axiom.
\item \( Q \) folgt aus der Informationserhaltung und dem Variationsprinzip.
\item Energie ist der Informationsfluss.
\item Masse ist der Widerstand gegen Informationsänderung (Amplitude): \( m = \nabla |I| \).
\item Ladung ist die eichinvariante Divergenz der Phase: \( q = \nabla^2 \phi \), diskret \( q_k = \sum_l (\phi_k - \phi_l) \).
\item Ladung ist eine Erhaltungsgröße – sie wird weder durch Fluktuation noch durch Energiesenke verändert.
\item \( Q \) ist das Bohm-Potential der DBT.
\item \( Q \) ist die Quelle der Gravitation (\( \vec{F} = -\nabla Q \)).
\item \( Q \) ist die Quelle der Elektrodynamik (\( \vec{F} = -\nabla Q \)).
\item \( Q \) ist das Prinzip der Selbstorganisation.
\end{itemize}

\( Q \) ist die Quelle der Physik.